\SourceLeaf{397}%
\noindent [Vorher wurde gesetzt:]
\[
\begin{aligned}
aa+bb+cc&=P\\
a\alpha+b\beta+c\gamma&=Q\\
\alpha\alpha+\beta\beta+\gamma\gamma&=R,
\end{aligned}
\]
[mithin wird:]
\[
\frac{\partial Q}{\partial t}=a'\alpha+b'\beta+c'\gamma+\tfrac12\frac{\partial P}{\partial\theta}.
\]

[Setzt man nun ein für allemal
\[
\begin{aligned}
\frac{\partial P}{\partial t}&=P',&\frac{\partial P}{\partial\theta}&=P'',\\
\frac{\partial Q}{\partial t}&=Q',&\frac{\partial Q}{\partial\theta}&=Q'',\\
\frac{\partial R}{\partial t}&=R',&\frac{\partial R}{\partial\theta}&=R'',
\end{aligned}
\]
so wird:
\[
\sigma s^3.\frac{\tau}{s}\sin\psi=P(Q'-\tfrac12P'')-\tfrac12QP'.
\]
Nun ist
\[
\frac{\tau}{s}\sin\psi.s\,du=\tau\sin\psi\,du
\]
das Differential der Seitenkrümmung, also ergibt sich, dass]
\[
(P(Q'-\tfrac12P'')-\tfrac12QP')dt
\]
gleich $\sigma ss$ mal dem Differential der Seitenkrümmung [ist. Es war aber]
\[
\begin{aligned}
\sigma\sigma&=PR-QQ\\
ss&=P;
\end{aligned}
\]
[demnach erhält man schliesslich die Formel:]
\begin{equation}\tag*{[14]}
\text{Diff. d. Seitenkr.}=\frac{(PQ'-\tfrac12PP''-\tfrac12QP')dt}{\BareRoot{(PR-QQ)}.P}.
\end{equation}
\Needspace{7\baselineskip}
\begin{center}\bfseries [4.]\end{center}
\noindent\hspace*{1.2em}Allgemein ist [das] Differential der Seitenkrümmung
\[
\begin{aligned}
df&=\tau\,du.\sin\psi\\
&=\tau\,du\{\xi'(Y\zeta-Z\eta)+\eta'(Z\xi-X\zeta)+\zeta'(X\eta-Y\xi)\},
\end{aligned}
\]
